%%%%%%%%%%  scan idx 232  |  folio 219  %%%%%%%%%%
$\operatorname{Ell}^{\underline n}(F)$ de
$\operatorname{Comp.diff}^{\underline n}(F)$ formé des complexes elliptiques d'ordre $\underline n$ est ouvert~; c'est cet espace (qui en général n'est pas localement compact) qui joue le rôle d'un espace de paramètres locaux universels.

\textbf{3.2.4.} Cependant, même si la généralisation de (3.2.1) à laquelle il a été fait allusion plus haut s'avère impossible, on peut toujours associer à $E^\bullet$, sous la seule hypothèse que la fibration $p$ soit propre et localement triviale, un élément de $K^\bullet(S)$ méritant le nom d'indice analytique de $E^\bullet$. En effet, supposons (pour simplifier\footnote{pour le cas général, utiliser la notion de «~fibré hilbertien faible~» ou champ continu d'espaces de Hilbert de [12].}) que l'on soit dans la situation «~locale~» où $X=S\times V$ ($V$, une variété compacte), $E^\bullet$ étant donné par une application continue de $S$ dans $\operatorname{Ell}^{\underline n}(F)$ (3.2.3). Les assertions (a') et (b) de la démonstration de (3.2.1) sont en fait valables sans aucune hypothèse sur $S$~: il est donc raisonnable de poser, avec les notations de (loc.cit.)~:
\[
  \operatorname{ind.an}(E^\bullet)=C\ell(\underline F_r^\bullet)\in K^\bullet(S),
\]
cet élément étant indépendant de $r$.

\textbf{3.2.5.} Supposons $S$ paracompact, notons $H$ un espace de Hilbert complexe de dimension hilbertienne infinie, $I(H)$ le sous-espace de $L(H,H)$ formé des opérateurs «~à indice~». Alors, les complexes $F_r^\bullet$ définissent [12] un élément canonique dans le groupe $[S,I(H)]$ des classes d'homotopie d'applications continues de $S$ dans $I(H)$, mais le rédacteur ignore pour $S$ non compact (même dans le cas où $S$ est localement compact) quelles sont les relations entre les groupes $[S,I(H)]$ et $K^\bullet(S)$ (existe-t-il une flèche naturelle de l'un vers l'autre~? sont-ils isomorphes~?), donc a fortiori entre les deux indices ainsi définis.

Pour une généralisation, sous les hypothèses de (3.2.1), du théorème de l'index d'Atiyah--Singer, nous renvoyons au travail récent de Shih [17], ou à un travail en préparation de M. F. Atiyah et I. M. Singer.

